\chapter[Fundamentação Teórica]{Fundamentação Teórica}\label{capitulo3-0}

A fundamentação teórica deste trabalho tem como objetivo apresentar conceitos, modelos e estudos que oferecem uma base científica para a compreensão tanto dos processos de avaliação psicopedagógica quanto do papel das tecnologias digitais no apoio às práticas educacionais. Ao reunir contribuições de áreas como desenvolvimento infantil, psicopedagogia, sistemas web e engenharia de software, este capítulo constrói o arcabouço necessário para identificar lacunas na literatura e justificar a relevância de uma solução tecnológica especificamente voltada à avaliação psicopedagógica na educação infantil.

\section{Desenvolvimento Infantil}

O desenvolvimento infantil envolve múltiplas dimensões cognitivas, motoras, sociais e emocionais, que se articulam de forma interdependente e constituem o núcleo de qualquer avaliação psicopedagógica consistente. Segundo \citeonline{Abreu_Barbosa_2025}, o desenvolvimento cognitivo, social e motor na infância está profundamente ligado a práticas pedagógicas lúdicas, à exploração de diferentes linguagens simbólicas e à mediação qualificada do brincar. Tais práticas ativam processos de raciocínio, linguagem, coordenação motora e interação social, contribuindo para trajetórias de aprendizagem mais ricas e diversificadas.

A dimensão emocional, frequentemente subestimada em abordagens estritamente conteudistas, é igualmente central. Estudos indicam que a afetividade influencia de maneira direta o desenvolvimento cognitivo, sobretudo na educação infantil, em que vínculos seguros e interações acolhedoras funcionam como base para a apropriação de novos conhecimentos \cite{Felix_Oliveira_Junior_2022}. Nessa perspectiva, a avaliação que desconsidera o contexto afetivo e relacional tende a produzir leituras empobrecidas das dificuldades e das potencialidades das crianças.

A produção clássica de Piaget, Vygotsky e Wallon constitui referência obrigatória para a compreensão dessas interações. Piaget destaca a construção ativa do conhecimento por meio de operações mentais progressivas, evidenciando que a criança organiza esquemas de ação e pensamento a partir da interação com o meio. Vygotsky enfatiza a mediação social, a linguagem e a zona de desenvolvimento proximal, sublinhando que o desenvolvimento ocorre em estreita relação com práticas culturais e interações com o outro. Wallon, por sua vez, confere centralidade à afetividade e ao corpo no desenvolvimento humano, mostrando a indissociabilidade entre emoção, movimento e cognição \cite{morgado2024_construcao_conhecimento}. Em conjunto, essas contribuições reforçam a ideia de que crianças pequenas são sujeitos ativos e emocionais, cujas trajetórias são moldadas por contextos sociais, simbólicos e afetivos.

A primeira infância, compreendida aproximadamente entre zero e seis anos, é reconhecida como período crítico para a organização cognitiva, emocional e social. O \citeonline{cpapi2021_primeira_infancia} ressalta que experiências vividas nessa faixa etária têm impacto duradouro sobre a estruturação neural e psicológica, o que implica responsabilidade acrescida nas práticas avaliativas e interventivas. Nesse cenário, a avaliação exige observação sistemática e instrumentos cuidadosamente escolhidos, capazes de captar dimensões quantitativas, como escores de linguagem ou coordenação motora, e dimensões qualitativas, como interações sociais, brincadeiras simbólicas e padrões de regulação emocional.

Um sistema digital de apoio à avaliação psicopedagógica precisa, portanto, estar orientado a registrar de forma estruturada essa multiplicidade de dimensões, favorecer o acompanhamento longitudinal e permitir devolutivas que expressem a singularidade de cada criança, em consonância com os referenciais teóricos aqui discutidos.

\section{Avaliação Psicopedagógica}

A avaliação psicopedagógica é um processo investigativo e integrador cujo objetivo é compreender condições e fatores que interferem na aprendizagem do sujeito, identificando dificuldades, potencialidades e contextos que as atravessam. Trata-se de uma prática multidimensional, que combina entrevistas, anamnese, observações, aplicação de instrumentos e análise interpretativa, culminando na elaboração de devolutivas que orientam famílias, escolas e outros profissionais. A literatura contemporânea insiste que se trata de um processo contextualizado e formativo e não apenas de uma operação de classificação diagnóstica \cite{pires2021_para_quem_serve}.

Nesse processo, buscam-se mapear processos cognitivos, afetivos e sociais relacionados à aprendizagem, organizando evidências que subsidiem intervenções pedagógicas e, quando necessário, encaminhamentos clínicos ou multidisciplinares. Os objetivos incluem identificar barreiras ao aprender, formular hipóteses explicativas de origem biológica, emocional, familiar ou institucional e propor estratégias de apoio. A natureza integradora e reflexiva da avaliação psicopedagógica supõe a articulação entre dados quantitativos, como escores de instrumentos padronizados, e dados qualitativos, como relatos, observações e produções espontâneas \cite{Cruz_Batista_de_Souza_dos_Santos_de_Lima_Amorim_de_Andrade_Pires_2023}.

A avaliação psicopedagógica pode assumir diferentes funções, a depender de sua intencionalidade. A avaliação diagnóstica tem por foco identificar a origem das dificuldades de aprendizagem, contemplando aspectos cognitivos, emocionais e sociais, em sintonia com as discussões de \citeonline{weiss2020_psicopedagogia_clinica}, \citeonline{bossa2023_psicopedagogia} e \citeonline{fernandez1991_inteligencia}. Na perspectiva clínica, contudo, diagnóstico e intervenção não constituem momentos estanques. \citeonline{visca1987_clinica_psicopedagogica} organiza a atuação psicopedagógica em diagnóstico, tratamento corretor e prevenção, e inicia o diagnóstico pela Entrevista Operativa Centrada na Aprendizagem (EOCA), em que o profissional propõe situações ao sujeito para observar sua modalidade de aprendizagem em ação. É nesse sentido que se fala em avaliação interventiva: as situações mediadas pelo profissional permitem observar o sujeito enquanto aprende, reavaliar hipóteses e ajustar estratégias ao longo do processo. A avaliação contínua corresponde ao acompanhamento sistemático e longitudinal, registrando avanços, recaídas e respostas às intervenções, sendo apontada como componente indispensável de práticas psicopedagógicas consistentes \cite{Santos_Silva_2020}. Essas perspectivas não se excluem; ao contrário, se articulam em um ciclo avaliativo que procura dar conta da complexidade da aprendizagem infantil.

\section{Etapas da Avaliação Psicopedagógica}

A literatura psicopedagógica apresenta um conjunto relativamente estável de etapas que compõem o processo de avaliação, ainda que autores e correntes adotem nomenclaturas diversas. Em linhas gerais, a avaliação psicopedagógica na infância compreende uma sequência de momentos articulados que, quando conduzidos de forma sistemática, favorecem a construção de hipóteses diagnósticas consistentes e o delineamento de intervenções contextualizadas. Entre essas etapas, podem ser destacadas:

\begin{enumerate}
    \item Entrevista inicial: a entrevista inicial constitui o ponto de partida, no qual se acolhe a queixa trazida pela família ou pela escola, se delimita o foco investigativo e se combinam aspectos logísticos do processo avaliativo. Estudos recentes ressaltam que uma escuta qualificada nesse momento permite situar a demanda em um contexto mais amplo, evitando conclusões precipitadas e orientando a escolha de instrumentos e estratégias de observação \cite{Franca_Oliveira_Oliveira_Ribeiro_Dias_Victoria_2024}.
    \item Anamnese: a anamnese consiste no levantamento detalhado de aspectos do desenvolvimento infantil, como marcos motores e de linguagem, histórico de saúde, sono, alimentação, percurso escolar, relações familiares e eventos significativos. Essa etapa é fundamental para diferenciar hipóteses de origem predominantemente pedagógica, socioemocional ou neurobiológica e para planejar, de forma mais precisa, os procedimentos subsequentes \cite{Transtorno_do_Espectro_Autista}.
    \item Observação psicopedagógica: a observação direta em contextos estruturados e espontâneos, como situações de brincadeira, sala de aula e interações com pares e adultos, permite captar comportamentos, estratégias de resolução de tarefas, modos de comunicação e respostas emocionais que nem sempre emergem em testes formais. Nessa etapa insere-se a EOCA \cite{visca1987_clinica_psicopedagogica}, na qual a criança é convidada a mostrar o que sabe fazer. A literatura recente enfatiza que, na educação infantil, a observação sistemática é recurso central para a formulação de hipóteses, pois revela aspectos qualitativos decisivos \cite{Costa_2024}.
    \item Aplicação de instrumentos e atividades avaliativas: esta etapa compreende o uso de tarefas lúdicas, provas operatórias, protocolos padronizados, checklists, escalas e produções espontâneas. Estudos sobre instrumentos psicopedagógicos apontam a importância de selecionar procedimentos validados e adequados à faixa etária e ao objetivo da avaliação e, quando adaptações forem necessárias, registrar de maneira explícita tais modificações, de modo a preservar a rastreabilidade das evidências \cite{tavares2022_instrumentos_psicopedagogicos}.
    \item Análise e integração de dados: após a coleta de informações, é necessário integrar dados oriundos da entrevista, da anamnese, das observações e dos instrumentos, confirmando, refinando ou descartando hipóteses. Pesquisas sobre práticas psicopedagógicas no Brasil destacam a relevância de considerar fatores contextuais relacionados à escola, à família e às condições socioeconômicas, evitando interpretações reducionistas e patologizantes \cite{Batista_Souza_Rosa_Nakano_Lima_2025}.
    \item Devolutiva: a devolutiva é o momento no qual se comunicam os achados à família e, quando pertinente, à escola, explicitando conclusões, fundamentos e encaminhamentos. A literatura ressalta que a devolutiva deve ser clara, ética e orientadora, estruturada de tal forma que as recomendações sejam compreendidas e factíveis, evitando jargões técnicos excessivos e rotulações estigmatizantes \cite{Franca_Oliveira_Oliveira_Ribeiro_Dias_Victoria_2024}.
    \item Plano de intervenção e acompanhamento: quando indicado, elabora-se um plano de intervenção ou de acompanhamento psicopedagógico, nas modalidades interventiva e contínua descritas na seção anterior. Estudos recentes defendem o uso de registros periódicos e sistemáticos para documentar progressos, rever estratégias e produzir evidências que sustentem decisões pedagógicas e clínicas \cite{Costa_2024, tavares2022_instrumentos_psicopedagogicos}.
\end{enumerate}

Essas etapas configuram um ciclo avaliativo altamente dependente de registros bem estruturados e de documentação consistente. A ausência de organização sistemática das informações tende a fragilizar o processo como um todo. Assim, ao propor um sistema web de apoio à avaliação psicopedagógica na educação infantil, torna-se indispensável traduzir esse fluxo em funcionalidades que permitam registrar, consultar e analisar, de modo integrado, o percurso de cada criança ao longo do processo avaliativo.

\section{Instrumentos e Ferramentas de Avaliação Psicopedagógica}\label{sec:fund_instrumentos}

No âmbito da psicopedagogia, emprega-se um conjunto diversificado de instrumentos com funções complementares, como avaliar o desenvolvimento do pensamento, detectar dificuldades de aprendizagem, compreender aspectos emocionais e sociais e orientar intervenções. Entre as ferramentas mais frequentes encontram-se testes com padrões definidos, tarefas inspiradas em provas operatórias, técnicas expressivas baseadas em desenhos e narrativas, registros de observação em forma de listas de verificação e questionários de comportamento voltados a indicadores como ansiedade, atenção e agitação \cite{Arimura_Scaglioni_Franco_de_Lima_2025}.

Estudos atuais sistematizam esses instrumentos, discutindo indicações, limites e formas de aplicação de acordo com o tipo de dificuldade, a faixa etária e o contexto de atendimento. Evidenciam também que a seleção de ferramentas não é neutra, mas deve estar alinhada às hipóteses formuladas a partir da entrevista inicial, da anamnese e das primeiras observações \cite{Arimura_Scaglioni_Franco_de_Lima_2025}. Nesse sentido, o uso de instrumentos sem ancoragem em um raciocínio clínico e pedagógico tende a produzir avaliações fragmentadas e pouco úteis.

Apesar da variedade de recursos disponíveis, grande parte da documentação produzida em contextos psicopedagógicos ainda ocorre em registros manuais dispersos, como anotações em papel, fichas isoladas e arquivos locais pouco integrados. Essa forma de registro traz implicações práticas e epistemológicas importantes. Do ponto de vista prático, dificulta a comparação longitudinal, aumenta o risco de perda de dados e limita a possibilidade de compartilhamento seguro entre profissionais. Do ponto de vista epistemológico, fragiliza a rastreabilidade das decisões, uma vez que adaptações, versões de instrumento e condições de aplicação nem sempre ficam registradas de maneira explícita \cite{Arimura_Scaglioni_Franco_de_Lima_2025}.

A padronização, a rastreabilidade e a contextualização dos dados emergem, assim, como requisitos centrais para a prática avaliativa contemporânea. Padronizar, neste contexto, não significa homogeneizar conclusões ou reduzir a complexidade dos casos, mas garantir que os registros incluam metadados mínimos, tais como identificador do profissional, data de aplicação, versão do instrumento, condições de aplicação e eventuais adaptações. Esses elementos permitem reanálises fundamentadas, comparações ao longo do tempo e maior transparência nas decisões.

A rastreabilidade é especialmente relevante quando instrumentos são adaptados ou combinados em protocolos específicos para determinado contexto educativo. Nessas situações, torna-se fundamental registrar quem realizou a adaptação, por que razões, em que data e quais campos foram alterados. Guias nacionais de avaliação e estudos sobre recomposição das aprendizagens enfatizam a importância de articular instrumentos padronizados e registros formativos contínuos, de modo a sustentar decisões pedagógicas em evidências consistentes \cite{mec2021_guia_avaliacao}.

Levantamentos sobre ferramentas de avaliação psicopedagógica indicam ainda grande heterogeneidade quanto ao público-alvo, ao formato de aplicação, aos critérios de interpretação e à qualidade dos manuais disponíveis \cite{tavares2022_instrumentos_psicopedagogicos}. Essa heterogeneidade reforça a necessidade de soluções tecnológicas capazes de organizar repositórios de instrumentos, vincular protocolos a metadados e facilitar a documentação sistemática dos atendimentos. Debates contemporâneos sobre avaliação também criticam abordagens centradas em patologias, sugerindo que os instrumentos devem registrar capacidades, contextos e situações, e não apenas déficits, o que demanda sistemas que comportem dados numéricos e descritivos de modo igualmente qualificado \cite{Barbosa_2025}.

Nesse cenário, a literatura aponta um conjunto de funcionalidades desejáveis para sistemas de apoio à prática psicopedagógica, entre as quais se destacam: repositório de instrumentos com metadados sobre versão e faixa etária; formulários eletrônicos para aplicação de protocolos e registro de observações qualitativas; mecanismos de controle de versões e registros de alteração; geração de relatórios em formatos interoperáveis, como PDF; perfis de acesso diferenciados para psicopedagogos, professores e famílias; e armazenamento seguro em conformidade com normas de privacidade e proteção de dados. Essas funcionalidades respondem diretamente às limitações dos registros manuais e às demandas de rastreabilidade, contextualização e segurança destacadas pela literatura.

\section{Sistemas Web na Educação}

Os sistemas web aplicados à educação configuram plataformas online que possibilitam a gestão, o registro, a análise e a comunicação de informações educacionais em tempo real e a partir de múltiplos perfis de usuário. Em contexto escolar, esses sistemas oferecem interfaces para professores, gestores, estudantes e famílias, centralizando dados como frequência, desempenho, atividades, trajetórias de aprendizagem e relatórios. O relatório TIC Educação 2023, do \citeonline{nic_tic_educacao2023}, indica que a digitalização desses dados contribui para maior transparência, consistência metodológica e apoio à tomada de decisão pedagógica baseada em evidências.

Entre as vantagens associadas a sistemas web, destaca-se a possibilidade de articular segurança da informação, escalabilidade e organização estruturada de dados. Conforme \citeonline{Silva_Sena_2025}, sistemas bem projetados devem garantir privacidade, integridade e controle de acesso, em consonância com a Lei Geral de Proteção de Dados. A escalabilidade permite que múltiplos usuários acessem simultaneamente sem comprometimento relevante de desempenho, o que é crucial para instituições com grande número de alunos e profissionais. O armazenamento estruturado facilita consultas, backups, auditorias e construção de relatórios, além de permitir a integração com outras bases de dados.

Apesar desses avanços, as plataformas educacionais de uso escolar raramente incluem módulos específicos para a prática psicopedagógica. Muitas delas concentram-se em funcionalidades administrativas, geração de boletins, controle de frequência e visualização de indicadores acadêmicos. Em geral, não contemplam espaços adequados para registro de observações clínicas, aplicação sistemática de instrumentos padronizados, documentação de anamnese ou elaboração de devolutivas psicopedagógicas. Além disso, a integração entre dados qualitativos e quantitativos é frequentemente limitada, o que dificulta uma visão integrada do desenvolvimento do aprendente.

A digitalização de dados educacionais sensíveis em sistemas web levanta ainda desafios técnicos e éticos. Estudos sobre plataformização da educação indicam riscos associados à centralização de informações em ambientes pouco transparentes quanto a políticas de uso, compartilhamento e armazenamento de dados \cite{internetlab2024_plataformizacao}. Pesquisas sobre adoção de plataformas e de inteligências artificiais na gestão educacional apontam lacunas na proteção da privacidade, demandando políticas institucionais claras e arquiteturas técnicas que privilegiem consentimento, minimização e criptografia de dados \cite{Borba_2025}. Estudos focados na implementação da Lei Geral de Proteção de Dados em instituições de ensino apontam desafios operacionais relevantes, como a necessidade de formação continuada, revisão contratual com fornecedores de tecnologia e classificação adequada dos dados sensíveis \cite{candiani2024_lgpd_instituicoes}. Levantamentos nacionais sobre o uso de tecnologias de informação e comunicação por crianças e adolescentes reforçam a urgência de práticas de segurança alinhadas às particularidades do público infantojuvenil \cite{cetic2024_tic_kids}.
\section{Fundamentos de Engenharia de Software}\label{sec:fundamentos_tecnicos}

Esta seção apresenta os conceitos de engenharia de software e de sistemas web que fundamentam as decisões técnicas do sistema desenvolvido, descritas no Capítulo~\ref{capitulo5}.

\subsection{Engenharia de requisitos e modelagem}\label{sec:fund_requisitos}

A engenharia de requisitos é o processo de descobrir, analisar, documentar e validar os serviços que um sistema deve oferecer e as restrições sob as quais ele deve operar \cite{sommerville2018_engenharia_software}. Os requisitos costumam ser classificados em funcionais, que descrevem os serviços e comportamentos esperados do sistema diante de determinadas entradas, e não funcionais, que estabelecem restrições e propriedades de qualidade, como segurança, desempenho, usabilidade e manutenibilidade. \citeonline{sommerville2018_engenharia_software} observa que requisitos não funcionais vagos são difíceis de verificar e recomenda, sempre que possível, expressá-los de forma mensurável. \citeonline{pressman2016_engenharia_software} destacam ainda que o levantamento de requisitos deve envolver os interessados no sistema, por meio de técnicas como entrevistas e análise de documentos, e que os requisitos devem ser validados antes e depois da implementação.

A modelagem complementa a especificação ao representar a estrutura e o comportamento do sistema em diferentes níveis de abstração. A \textit{Unified Modeling Language} (UML) padroniza diagramas para esse fim, como os de casos de uso, de classes e de sequência \cite{omg2017_uml}. Para a estrutura dos dados, o modelo entidade-relacionamento, proposto por \citeonline{chen1976_entity_relationship}, representa as entidades do domínio, seus atributos e os relacionamentos entre elas, incluindo a cardinalidade de cada relacionamento, e é amplamente empregado no projeto conceitual de bancos de dados relacionais.

\subsection{Arquitetura em camadas e REST}\label{sec:fund_arquitetura}

A arquitetura em camadas organiza um sistema em níveis com responsabilidades distintas, em que cada camada utiliza apenas os serviços da camada imediatamente inferior. A separação mais comum em aplicações de informação distingue as camadas de apresentação, que trata da interação com o usuário, de domínio ou aplicação, que concentra as regras de negócio, e de fonte de dados, que trata da persistência \cite{fowler2002_patterns}. Essa organização favorece a manutenção, pois alterações em uma camada tendem a não se propagar às demais, e permite substituir ou testar uma camada isoladamente \cite{sommerville2018_engenharia_software}.

Em aplicações web, a comunicação entre a interface executada no navegador e o servidor costuma seguir o estilo arquitetural REST (\textit{Representational State Transfer}), definido por \citeonline{fielding2000_rest}. O REST organiza o sistema em torno de recursos identificados por URIs e manipulados por meio de uma interface uniforme, e impõe restrições como a separação entre cliente e servidor e a ausência de estado de sessão no servidor: cada requisição deve conter todas as informações necessárias para ser processada. Na prática, as APIs REST utilizam os métodos do protocolo HTTP (GET, POST, PUT, PATCH e DELETE) para expressar as operações sobre os recursos e os códigos de estado da resposta para indicar o resultado, como 401 para requisições sem autenticação válida, 403 para acesso não autorizado e 429 para excesso de requisições \cite{rfc9110_http}.

\subsection{Bancos de dados relacionais e mapeamento objeto-relacional}\label{sec:fund_banco}

O modelo relacional, proposto por \citeonline{codd1970_relational}, representa os dados como relações, ou tabelas, compostas por tuplas e atributos, e manipula-os por meio de operações declarativas, independentes da forma de armazenamento físico. A integridade dos dados é assegurada por restrições como as chaves primárias, que identificam cada tupla de forma única, e as chaves estrangeiras, que garantem que um registro só faça referência a registros existentes em outra relação. Os sistemas gerenciadores de bancos de dados relacionais oferecem ainda transações, que permitem executar um conjunto de operações de forma atômica.

Como as linguagens de programação orientadas a objetos e o modelo relacional representam os dados de maneiras diferentes, é comum o uso de uma camada de mapeamento objeto-relacional (\textit{Object-Relational Mapping}, ORM), responsável por converter os registros das tabelas em objetos da aplicação e vice-versa. \citeonline{fowler2002_patterns} descreve os padrões que fundamentam essas ferramentas, como o \textit{Data Mapper}, que mantém os objetos de domínio independentes do banco de dados. O uso de um ORM reduz o código repetitivo de acesso a dados e, quando as consultas são construídas por meio da própria ferramenta, evita a concatenação manual de comandos SQL, uma das causas de vulnerabilidades de injeção.

\subsection{Autenticação com JWT e controle de acesso baseado em papéis}\label{sec:fund_seguranca}

A autenticação verifica a identidade de quem acessa o sistema, enquanto a autorização determina o que cada usuário autenticado pode fazer. Em APIs REST, a ausência de estado de sessão no servidor favorece o uso de credenciais autocontidas, como o JSON Web Token (JWT), definido pela RFC 7519 \cite{rfc7519_jwt}. Um JWT é composto por um cabeçalho, um conjunto de declarações (\textit{claims}), como o identificador do usuário e o momento de expiração, e uma assinatura que permite ao servidor verificar que o conteúdo não foi alterado. Como o token assinado permanece válido até expirar, é comum combinar um token de acesso de curta duração com um token de renovação (\textit{refresh token}) de duração maior, que pode ser revogado no servidor. As senhas, por sua vez, não devem ser armazenadas em texto plano, mas como resultado de funções de \textit{hash} deliberadamente lentas, como o bcrypt, cujo fator de custo ajustável permite aumentar o esforço computacional necessário para ataques de força bruta à medida que o poder de processamento cresce \cite{provos1999_bcrypt}.

No controle de acesso baseado em papéis (\textit{Role-Based Access Control}, RBAC), as permissões não são atribuídas diretamente aos usuários, mas a papéis que representam funções na organização; os usuários recebem papéis e, por meio deles, as permissões correspondentes \cite{sandhu1996_rbac}. Esse modelo simplifica a administração das permissões, mas não expressa, por si só, restrições que dependem da relação entre o usuário e um registro específico, como a regra de que um profissional só pode acessar os dados das crianças que acompanha. Restrições desse tipo são tratadas pelo controle de acesso baseado em atributos (\textit{Attribute-Based Access Control}, ABAC), em que a decisão considera atributos do usuário, do recurso e do contexto, como o responsável por um registro \cite{hu2014_abac}.

Como referência para a segurança de aplicações web, a OWASP mantém o OWASP Top 10, que reúne as categorias de risco mais críticas para essas aplicações, entre elas as falhas de controle de acesso, as falhas criptográficas e as falhas de identificação e autenticação \cite{owasp2021_top10}. De forma complementar, o \textit{Application Security Verification Standard} (ASVS) organiza requisitos de segurança verificáveis, agrupados em níveis de exigência, que podem orientar tanto o desenvolvimento quanto a verificação de controles como autenticação, gerenciamento de sessão e controle de acesso \cite{owasp2025_asvs}.

\subsection{Testes de software}\label{sec:fund_testes}

O teste de software é o processo de executar um programa com a intenção de encontrar defeitos \cite{myers2011_art_testing}. \citeonline{sommerville2018_engenharia_software} distingue os testes de desenvolvimento em níveis: os testes unitários verificam unidades isoladas, como funções ou classes; os testes de componente ou de integração verificam a interação entre unidades por meio de suas interfaces; e os testes de sistema verificam o sistema completo. Nos testes unitários e de integração, dependências externas, como o banco de dados, podem ser substituídas por objetos simulados, o que torna os testes rápidos e determinísticos, mas exige verificações complementares contra as dependências reais. A cobertura de código, que mede a proporção do código executada pelos testes, é um indicador do quanto o programa foi exercitado, mas não garante a ausência de defeitos \cite{myers2011_art_testing}. Os testes automatizados distinguem-se ainda da avaliação com usuários, que verifica se o sistema atende às necessidades de quem o utiliza, aspecto que os testes técnicos não alcançam \cite{sommerville2018_engenharia_software}.

\section{Requisitos para Sistemas de Avaliação e Registros Clínicos}

A construção de sistemas digitais destinados à avaliação psicopedagógica envolve a observância de requisitos técnicos, legais e de usabilidade que assegurem confiança, utilidade e aderência às práticas profissionais. Esses requisitos abrangem padrões de usabilidade e acessibilidade, segurança da informação, em especial no que se refere à Lei Geral de Proteção de Dados, e funcionalidades específicas para registro longitudinal, logs de auditoria e controle de versões, além da interoperabilidade com outros sistemas.

Sistemas destinados à área educacional e psicopedagógica devem priorizar usabilidade e acessibilidade. Isso implica interfaces claras, organização lógica das tarefas, rotinas consistentes de navegação e atenção às diretrizes internacionais de acessibilidade, como as recomendações das Web Content Accessibility Guidelines \cite{w3c_wcag22_pt}. Guias nacionais ressaltam a importância de incorporar critérios de acessibilidade desde a fase de concepção, com uso de listas de verificação e testes com públicos diversos \cite{govbr_guia_acessibilidade}. Em contextos clínicos e educacionais, em que o tempo de atendimento é limitado e o volume de demandas é elevado, a experiência de uso tem impacto direto sobre a qualidade do registro.

A interoperabilidade torna-se relevante quando se busca integrar o registro psicopedagógico a outros sistemas escolares ou de saúde, como prontuários ou plataformas de redes municipais. Estudos sobre o tema mostram que modelos de dados padronizados e interfaces de programação bem definidas facilitam trocas seguras de informação e reduzem redundâncias \cite{Lins_Ferreira_Nutels_Junior_Tenorio_Filho_2024, mota2024_rnds_interoperabilidade}. Por depender de padrões de dados ainda não consolidados para a psicopedagogia e de acordos com os sistemas parceiros, a interoperabilidade não foi incluída no escopo deste trabalho, que se concentra no registro e na organização do processo avaliativo dentro do próprio sistema; ela é retomada como trabalho futuro no Capítulo~\ref{capitulo7}.

A confiabilidade das avaliações psicopedagógicas depende também da presença de recursos para registro longitudinal, logs de auditoria e controle de versões. O registro longitudinal permite acompanhar a trajetória de cada criança, conservando o histórico de atendimentos, resultados de instrumentos, observações e devolutivas. Logs de auditoria registram acessos e alterações de dados, indicando quem realizou a ação, em que momento e com que conteúdo, o que é fundamental para rastreabilidade e responsabilização. O controle de versões permite documentar adaptações em protocolos, garantindo que comparações temporais considerem mudanças em instrumentos e procedimentos.

Além da arquitetura modular em camadas (Seção~\ref{sec:fund_arquitetura}), a literatura técnica recomenda a realização de testes iterativos de usabilidade, planos de governança de dados, políticas de anonimização quando apropriadas para fins de pesquisa e a participação de diferentes grupos de interesse desde o levantamento de requisitos até a validação de protótipos \cite{Costa_Camargos_Viana_Mendes_2025}. A conformidade com a legislação sobre proteção de dados demanda, por sua vez, controles de acesso, mecanismos de consentimento, criptografia e políticas claras de retenção e descarte de dados.

\section{Soluções Existentes no Mercado}\label{sec:solucoes_mercado}

Para verificar se a lacuna apontada na literatura se mantém diante das soluções comerciais disponíveis, realizou-se um levantamento exploratório de sistemas de prontuário eletrônico oferecidos no mercado brasileiro a psicopedagogos, psicólogos e clínicas multiprofissionais. A análise baseou-se nas funcionalidades divulgadas publicamente nos sítios oficiais de cada produto, consultados em novembro de 2026; portanto, a ausência de um recurso no Quadro~\ref{qua:solucoes_mercado} indica que ele não é informado pelo fornecedor, e não necessariamente que inexista no produto.

O iPsy Pro é um sistema de gestão voltado a psicopedagogos, que organiza agenda, financeiro e um prontuário estruturado em resumo, anamnese, avaliação, plano e sessões, além de oferecer modelos prontos de documentos e um banco de jogos para uso em sessão. O próprio fornecedor esclarece que a ferramenta não substitui a avaliação nem os instrumentos de uso restrito, de modo que não há aplicação estruturada de protocolos no sistema \cite{ipsy_pro}. O Cerebraris, também destinado a psicopedagogos e neuropsicopedagogos, avança nesse aspecto ao permitir a aplicação e o registro de um conjunto predefinido de instrumentos, como TDE-II, SNAP-IV e M-CHAT, com pontuação automática, além de anamnese, evolução e geração de laudos a partir de modelos prontos \cite{cerebraris}. Não são informados, contudo, recursos para que o profissional cadastre instrumentos próprios ou adaptados.

Entre as soluções generalistas, o PsicoManager, voltado a psicólogos e clínicas, reúne prontuário eletrônico, documentos, escalas, agenda, financeiro e gestão de múltiplos profissionais com permissões por usuário \cite{psicomanager}. O Ninsaúde Clinic, sistema multiprofissional, permite personalizar protocolos e formulários do prontuário a partir de uma biblioteca por especialidade ou de modelos criados pelo próprio usuário \cite{ninsaude_prontuario}. Nesses casos, entretanto, os formulários funcionam como modelos de prontuário, e não há informação sobre como alterações em um modelo afetam os registros já preenchidos com ele.

\begin{quadro}[H]
\centering
\small
\caption{Funcionalidades divulgadas por soluções de prontuário eletrônico no mercado brasileiro}
\begin{tabular}{p{2.2cm} p{2.3cm} p{2.6cm} p{2.1cm} p{2.4cm} p{1.6cm}}
\hline
\textbf{Produto} & \textbf{Público} & \textbf{Instrumentos de avaliação} & \textbf{Criação de instrumentos próprios} & \textbf{Versionamento vinculado às aplicações} & \textbf{Agenda e financeiro} \\
\hline
iPsy Pro & Psicopedagogos & Não (declarado pelo fornecedor) & Não informado & Não informado & Sim \\
Cerebraris & Psicopedagogos e neuropsicopedagogos & Conjunto predefinido, com pontuação automática & Não informado & Não informado & Sim \\
PsicoManager & Psicólogos e clínicas & Escalas (sem detalhamento) & Não informado & Não informado & Sim \\
Ninsaúde Clinic & Clínicas multiprofissionais & Formulários de prontuário & Sim (modelos de formulário) & Não informado & Sim \\
Sistema proposto & Psicopedagogos & Cadastrados pelo profissional & Sim & Sim & Não \\
\hline
\end{tabular}
\fonte{Elaborado pela autora (2026), com base em \citeonline{ipsy_pro}, \citeonline{cerebraris}, \citeonline{psicomanager} e \citeonline{ninsaude_prontuario}.}
\label{qua:solucoes_mercado}
\end{quadro}

O levantamento mostra que a existência de sistemas específicos para psicopedagogos é uma realidade e que esses produtos já atendem bem às rotinas administrativas do consultório, como agenda, cobrança e emissão de documentos, funcionalidades que não fazem parte do escopo deste trabalho. A lacuna identificada é mais restrita e diz respeito ao tratamento do instrumento avaliativo como um objeto de primeira classe no sistema: nenhuma das soluções analisadas divulga, simultaneamente, a possibilidade de o profissional cadastrar e adaptar seus próprios instrumentos e o controle de versões que preserve, em cada aplicação, a versão do instrumento efetivamente utilizada. Esse controle é o que permite comparar aplicações realizadas em momentos distintos sabendo exatamente quais questões foram respondidas em cada uma, requisito discutido na seção anterior.

\section{Síntese Integradora: Lacunas e Oportunidades}

A literatura recente sobre tecnologias educacionais e psicopedagogia evidencia o potencial das plataformas digitais para registrar evidências, apoiar a tomada de decisão e ampliar o alcance das práticas avaliativas, especialmente em um contexto de intensificação do uso de recursos digitais. Esses estudos ressaltam que o êxito de tais plataformas depende da aderência às especificidades da infância, da formação adequada dos profissionais e da existência de políticas de uso que preservem direitos e garantias dos sujeitos envolvidos \cite{Aielo_Silva_Augusto_2025}.

Ao mesmo tempo, persistem as lacunas discutidas nas seções anteriores: as plataformas educacionais generalistas não incluem módulos dedicados à prática psicopedagógica, e permanece a dificuldade de articular, em um mesmo ambiente, dados quantitativos e qualitativos sem perder validade interpretativa. Adicionalmente, pesquisas sobre plataformização da educação e sobre infraestrutura tecnológica apontam obstáculos práticos, como conectividade restrita, desigualdade de acesso e insuficiente capacitação de docentes e equipes técnicas \cite{edu_plataformas2024}. Tais desafios repercutem na adoção real de soluções digitais, que muitas vezes permanecem subutilizadas ou são incorporadas de modo fragmentado. Nesse cenário, propostas tecnológicas mais específicas, concebidas com base em diagnósticos precisos das necessidades de um campo profissional, podem contribuir para superar parte dessas limitações.

A revisão realizada indica, portanto, oportunidades concretas para um sistema web voltado à avaliação psicopedagógica na educação infantil, que correspondem às funcionalidades apontadas na Seção~\ref{sec:fund_instrumentos}: repositório de instrumentos com controle de versões, registro estruturado de aplicações e observações qualitativas, relatórios interpretativos e requisitos de segurança, privacidade e acessibilidade incorporados desde a concepção \cite{Almeida_2024}.
